\documentclass[12pt,a4paper]{article}
\usepackage[utf8]{inputenc}
\usepackage[T1]{fontenc}
\usepackage[brazil]{babel}
\usepackage[a4paper,margin=2.5cm]{geometry}
\usepackage{graphicx}
\usepackage{listings}
\usepackage{xcolor}
\usepackage{hyperref}
\usepackage{indentfirst}
\usepackage{titlesec}

% Configuração do estilo de código
\definecolor{codegray}{rgb}{0.5,0.5,0.5}
\definecolor{codepurple}{rgb}{0.58,0,0.82}
\definecolor{backcolour}{rgb}{0.95,0.95,0.92}
\definecolor{codegreen}{rgb}{0,0.6,0}

\lstdefinestyle{cstyle}{
    backgroundcolor=\color{backcolour},
    commentstyle=\color{codegreen},
    keywordstyle=\color{blue},
    numberstyle=\tiny\color{codegray},
    stringstyle=\color{codepurple},
    basicstyle=\ttfamily\footnotesize,
    breakatwhitespace=false,
    breaklines=true,
    captionpos=b,
    keepspaces=true,
    numbers=left,
    numbersep=5pt,
    showspaces=false,
    showstringspaces=false,
    showtabs=false,
    tabsize=2,
    language=C
}

\lstset{style=cstyle}

\begin{document}

% -------------------------------------------------------
% CAPA
% -------------------------------------------------------
\begin{titlepage}
    \centering
    {\Large Universidade Federal de Mato Grosso do Sul\par}
    \vspace{0.8cm}
    \includegraphics[width=5cm]{logo_ufms.png}
    \vspace{0.8cm}\\
    {\large Faculdade de Computação\par}
    {\large Estruturas de Dados\par}
    {\large Professor: Anderson Bessa da Costa\par}
    \vfill
    {\Huge\bfseries Trabalho Prático -- Árvore B\par}
    \vfill
    \large
    Otávio Gabriel de Oliveira Santana -- 202519050583\\
    Guilherme Duarte Balduino -- 202519050680\\
    \vfill
    Campo Grande -- MS\\
    Maio/Junho 2026
\end{titlepage}

% -------------------------------------------------------
% SUMÁRIO
% -------------------------------------------------------
\tableofcontents
\newpage

% -------------------------------------------------------
% 1. INTRODUÇÃO
% -------------------------------------------------------
\section{Introdução}

As árvores B são árvores balanceadas utilizadas para facilitar e minimizar o acesso a memórias secundárias. As árvores B possuem diversas características, entre elas: múltiplas chaves por nó; alta ordem (quanto maior a sua ordem, mais chaves por nó, minimizando o número de acessos); e balanceamento garantido — todas as folhas estão sempre no mesmo nível, independentemente da ordem de inserção.

As árvores B possuem uma série de regras que devem ser seguidas:

\begin{enumerate}
    \item[(i)] a raiz é uma folha ou possui no mínimo dois filhos;
    \item[(ii)] cada nó interno (exceto a raiz) possui no mínimo $d + 1$ filhos;
    \item[(iii)] cada nó possui no máximo $2d + 1$ filhos;
    \item[(iv)] todas as folhas estão no mesmo nível.
\end{enumerate}

Neste trabalho tivemos como objetivo montar uma árvore B com ordem $D \geq 2$, aplicando duas funções principais (Busca e Inserção) e algumas subfunções auxiliares. O objetivo final era construir uma árvore B funcional que funcionasse como uma espécie de \textit{pokédex}, utilizando como material de estudos o livro \textit{Estruturas de Dados e seus Algoritmos} de Jayme Szwarcfiter e Liane Markenzon.

% -------------------------------------------------------
% 2. INSTRUÇÕES
% -------------------------------------------------------
\section{Instruções}

Foi instruído que seguíssemos as lógicas e algoritmos apresentados por Jayme como base para os estudos na implementação de uma árvore B. Teríamos acesso às bibliotecas padrões da linguagem C (ANSI C).

Há de ser instituída uma implementação do tipo \texttt{struct} que representasse os nós da árvore, onde cada nó deveria armazenar um conjunto de chaves (\textit{strings}) e ponteiros para seus filhos, conforme a definição da estrutura. A ordem da árvore deveria ser fornecida pelo usuário no início do programa, como forma de torná-la mais interativa (sendo $D \geq 2$).

Foi entregue pelo professor um arquivo contendo o nome de Pokémons que deveriam ser inseridos na árvore de maneira automática. Após a inserção dos nomes, o menu deveria ser apresentado, dando ao usuário as seguintes opções de navegação:

\begin{lstlisting}[language={}, numbers=none, backgroundcolor=\color{backcolour}]
// ----- // ----- // ARVORE B // ----- // ----- //
[1] - Buscar
[2] - Inserir
[9] - Finalizar
-------------------------------
Entre com a sua opcao:
\end{lstlisting}

Com o seguinte comportamento:

\begin{itemize}
    \item \textbf{Buscar:} lê uma chave \texttt{x} e realiza sua busca na árvore B, informando se ela foi encontrada e, quando aplicável, a posição onde se encontra;
    \item \textbf{Inserir:} lê uma chave \texttt{x} e realiza sua inserção na árvore B;
    \item \textbf{Finalizar:} encerra a execução do programa.
\end{itemize}

Sem a necessidade de implementar uma função de remoção.

% -------------------------------------------------------
% 3. IMPLEMENTAÇÃO
% -------------------------------------------------------
\section{Implementação}

\subsection{Struct}

De início foi implementado o \texttt{struct no}, que serviria de base para qualquer operação que iríamos realizar:

\begin{lstlisting}
typedef struct no
{
    int n;                     // Quantidade de chaves no No
    char *chaves[2*D_MAX];     // Lista de Strings -> Cada chave e uma String
    struct no* filhos[2*D_MAX+1];
} No;

No* raiz = NULL; // Definicao da raiz principal utilizada nas operacoes
int D;
\end{lstlisting}

Escolhemos utilizar ponteiros para strings alocadas dinamicamente com \texttt{strdup}, permitindo que $D$ seja dinâmico. Cada \texttt{chaves[i]} é um \texttt{char*} que aponta para uma string alocada independentemente na memória, evitando que uma string sobrescreva outra.

\subsection{Menu}

Sendo seu retorno um \texttt{int} (1, 2, 9 ou 0), a função \texttt{menu()} foi definida como a única interface visível ao usuário, com três opções possíveis: 1 para Busca, 2 para Inserção e 9 para encerrar o programa. Caso nenhuma opção válida seja selecionada, o programa informa o erro e retorna 0.

\begin{lstlisting}
while (1) {
    if (menu() == 9) { // Caso o usuario selecione [9] - Finalizar
        break;
    }
}
\end{lstlisting}

A função \texttt{menu()} é chamada na \texttt{main} dentro de um loop que só é quebrado quando o usuário digita 9.

\subsection{Busca}

Realizada de forma recursiva, a função \texttt{busca()} determina se um elemento está na estrutura ou não. O caso base é verificado checando se o ponteiro \texttt{No*} é nulo — caso seja, o elemento não está na árvore:

\begin{lstlisting}
int busca(No* pt, char *x) {
    if (pt == NULL) {
        printf("Pokemon '%s' nao encontrado :(\n", x);
        return 0;
    }
\end{lstlisting}

Caso contrário, percorremos as chaves do nó da esquerda para a direita utilizando uma dupla condição no \texttt{while}:

\begin{lstlisting}
int i = 0;
// Vai percorrendo as chaves (Esq->Dir)
while (i < pt->n && strcmp(x, pt->chaves[i]) > 0) {
    i++;
}
\end{lstlisting}

A função \texttt{strcmp}, nativa da biblioteca \texttt{<string.h>}, faz a comparação entre duas strings e retorna, sendo \texttt{Y} o valor retornado:

\begin{itemize}
    \item \texttt{Y > 0} — caso o primeiro parâmetro seja alfabeticamente maior que o segundo;
    \item \texttt{Y == 0} — caso sejam iguais;
    \item \texttt{Y < 0} — caso o primeiro parâmetro seja alfabeticamente menor que o segundo.
\end{itemize}

Esta dupla condição dentro do \texttt{while} garante que \texttt{i} não ultrapasse o limite do nó e que avance enquanto a chave buscada for maior. Ao sair do loop, há dois possíveis cenários:

\textbf{Caso 1 — \texttt{strcmp} retorna 0:} a string foi encontrada dentro do nó.

\begin{lstlisting}
if (i < pt->n && strcmp(x, pt->chaves[i]) == 0) {
    printf("Pokemon '%s' encontrado! Posicao %d de seu No :)\n", x, i);
    return 1;
}
\end{lstlisting}

\textbf{Caso 2 — \texttt{i} ultrapassou seu valor máximo:} o elemento não está neste nó. A busca desce recursivamente por \texttt{filhos[i]}, que é o único filho onde o elemento poderia estar:

\begin{lstlisting}
return busca(pt->filhos[i], x); // Desce para o filho[i]
\end{lstlisting}

\subsection{Add\_Key}

A função \texttt{add\_key()} é utilizada tanto na função \texttt{insere()} quanto na função \texttt{split()}, sendo responsável por inserir um novo elemento dentro de uma página já existente, mantendo a integridade e a organização alfabética da árvore.

A função percorre as chaves do nó da esquerda para a direita comparando alfabeticamente com a nova chave. Ao encontrar a posição correta, ela desloca as demais chaves à sua direita uma posição à frente e insere o elemento no local correto:

\begin{lstlisting}
for (int j = nmr; j > i; j--) {
    pt->chaves[j] = pt->chaves[j-1];
}
pt->chaves[i] = key;
pt->n++;
\end{lstlisting}

\subsection{Split}

A função \texttt{split()} é acionada quando há sobrecarga em uma página, ou seja, quando o número de chaves atinge o máximo permitido ($2 \times D$). Ela divide a página em duas — metade esquerda e metade direita — e sobe o elemento do meio para o nó pai, que se torna o pai dessas duas subdivisões.\\

Optou-se pela estratégia de divisão preventiva (\textit{top-down}): antes de
descer para um filho, verifica-se se ele está cheio e, em caso afirmativo,
ele é dividido. Isso garante que a inserção em uma folha nunca encontre um
nó cheio, simplificando o tratamento recursivo.\\

Em um nó cheio com $2D$ chaves (índices $0$ a $2D-1$), o elemento do meio
escolhido para subir ao pai é o de índice \texttt{meio = D}. Assim, a metade
esquerda permanece com $D$ chaves (índices $0$ a $D-1$) no nó original, o
elemento de índice $D$ sobe, e a metade direita (índices $D+1$ a $2D-1$)
forma o novo nó.

Para realizar a divisão, é criado um novo nó \texttt{no\_dir} que recebe as chaves a partir do centro da página original. As chaves da metade direita são copiadas para \texttt{no\_dir} e seus espaços originais são transformados em \texttt{NULL}, limpando o nó original:

\begin{lstlisting}
while (j < 2*D) {
    no_dir->chaves[k] = pt->chaves[j];
    pt->chaves[j] = NULL; // Limpa o elemento do no original
    no_dir->n++;
    j++; k++;
}
\end{lstlisting}

O nó original \texttt{pt} passa a representar apenas a metade esquerda, com \texttt{pt->n = meio}.

Caso o nó dividido \textbf{não seja uma folha}, os filhos da metade direita também precisam ser transferidos para o \texttt{no\_dir}, garantindo que a árvore não perca os ponteiros para esses nós:

\begin{lstlisting}
if (pt->filhos[0] != NULL) {
    while (f <= 2*D) {
        no_dir->filhos[g] = pt->filhos[f];
        pt->filhos[f] = NULL;
        f++; g++;
    }
}
\end{lstlisting}

Por fim, cuida-se da redistribuição dos ponteiros dos filhos. Caso o nó dividido \textbf{não seja a raiz} (\texttt{posi != -1}), o \texttt{no\_dir} é encaixado como filho do pai na posição correta, empurrando os filhos subsequentes uma posição para a direita:

\begin{lstlisting}
for (int k = pai->n; k > posi+1; k--) {
    pai->filhos[k] = pai->filhos[k-1];
}
pai->filhos[posi+1] = no_dir;
\end{lstlisting}

Caso a página sobrecarregada seja a própria raiz (\texttt{posi == -1}), o elemento do meio se torna uma nova raiz e as duas metades tornam-se seus filhos:

\begin{lstlisting}
raiz = novo_no(0);
raiz->chaves[0] = chave_meio;
raiz->n++;
raiz->filhos[0] = pt;
raiz->filhos[1] = no_dir;
\end{lstlisting}

\subsection{Inserção}

A função \texttt{insere()} é responsável por inserir um novo elemento na árvore B, garantindo que todas as propriedades da estrutura sejam mantidas após a inserção. Ela trata as seguintes condições:

\textbf{Condição 1 — Árvore vazia:} caso a raiz seja \texttt{NULL}, um novo nó é criado e a chave é inserida diretamente nele.

\begin{lstlisting}
if (raiz == NULL) {
    raiz = novo_no(0);
    raiz->chaves[0] = nome;
    raiz->n++;
    return;
}
\end{lstlisting}

\textbf{Condição 2 — Raiz cheia:} caso a raiz esteja cheia antes de descer, é realizado um split preventivo para garantir que sempre haja espaço ao descer.

\begin{lstlisting}
if (pai == NULL && pt->n >= 2*D) {
    split(NULL, pt);
    pt = raiz;
}
\end{lstlisting}

\textbf{Condição 3 — Elemento já existe:} caso o elemento já esteja na árvore, nenhuma ação é realizada, evitando duplicatas.

\begin{lstlisting}
if (i < pt->n && strcmp(nome, pt->chaves[i]) == 0) {
    printf("Elemento ja se encontra na arvore.\n");
    return;
}
\end{lstlisting}

\textbf{Condição 4 — Chegou na folha:} \texttt{filhos[i] == NULL} indica que estamos numa folha. A função \texttt{add\_key()} é acionada para inserir a chave na posição correta.

\begin{lstlisting}
else if (pt->filhos[i] == NULL) {
    add_key(pt, nome);
    return;
}
\end{lstlisting}

\textbf{Condição 5 — Filho cheio:} caso o filho para onde desceríamos esteja cheio, é realizado um split preventivo nele antes de descer. Como o split reorganiza as chaves do pai, o valor de \texttt{i} é recalculado antes de descer recursivamente.

\begin{lstlisting}
else if (pt->filhos[i]->n >= 2*D) {
    split(pt, pt->filhos[i]);
    i = 0; // Recalcula i pois o pai foi reorganizado
    while (i < pt->n && strcmp(nome, pt->chaves[i]) > 0) {
        i++;
    }
    return insere(pt, pt->filhos[i], nome);
}
\end{lstlisting}

\textbf{Condição 6 — Filho não cheio:} caso o filho não esteja cheio, a função desce recursivamente por ele.

\begin{lstlisting}
else {
    return insere(pt, pt->filhos[i], nome);
}
\end{lstlisting}

% -------------------------------------------------------
% 4. DIFICULDADES ENCONTRADAS
% -------------------------------------------------------
\section{Dificuldades Encontradas}

\textbf{Split:} A principal dificuldade na função \texttt{split()} foi compreender que a divisão não deveria ocorrer apenas no filho, mas também tratar o pai quando ele fosse exceder o limite de chaves. Inicialmente estava sendo realizado apenas o split do filho, ignorando a propagação para o nó pai.\\

\textbf{Insere:} Na função \texttt{insere()} encontramos um problema com a primeira chamada sendo realizada com o pai \texttt{NULL} — sempre que um novo Pokémon era inserido, o anterior era sobrescrito pelo novo. O problema foi resolvido tratando corretamente o caso da raiz e garantindo que cada string fosse armazenada de forma independente com \texttt{strdup}.\\

\textbf{Ordenação das chaves:} Manter as chaves ordenadas dentro de cada nó foi um desafio, pois diversas situações podiam ocorrer ao inserir um novo elemento — no início, no meio ou no final. Para resolver isso, foi criada a função auxiliar\\
\texttt{add\_key()}, que centraliza e trata todas essas situações.\\

\textbf{Busca:} A dificuldade foi visualizar como percorrer os nós em ordem alfabética. O uso da função \texttt{strcmp()} facilitou muito esse processo, permitindo comparar as chaves de forma simples e decidir se deveríamos avançar para a direita ou descer para um filho.

% -------------------------------------------------------
% 5. FUNCIONALIDADES NÃO IMPLEMENTADAS
% -------------------------------------------------------
\section{Funcionalidades Não Implementadas}

Conforme especificado no enunciado do trabalho, a operação de \textbf{remoção} não foi implementada.

% -------------------------------------------------------
% 6. AMBIENTE DE DESENVOLVIMENTO
% -------------------------------------------------------
\section{Ambiente de Desenvolvimento}

O trabalho foi desenvolvido utilizando o seguinte ambiente:

\begin{itemize}
    \item \textbf{Sistema Operacional:} Windows 11
    \item \textbf{Editor:} Visual Studio Code
    \item \textbf{Compilador:} GCC versão 14.2.0
\end{itemize}

% -------------------------------------------------------
% 7. INSTRUÇÕES DE COMPILAÇÃO E EXECUÇÃO
% -------------------------------------------------------
\section{Instruções de Compilação e Execução}

Para compilar o programa, execute o seguinte comando no terminal dentro do diretório do projeto:

\begin{lstlisting}[language=bash, numbers=none]
gcc main_b-tree.c -o btree.exe
\end{lstlisting}

Para executar o programa:

\begin{lstlisting}[language=bash, numbers=none]
./btree.exe
\end{lstlisting}

\textbf{Observação:} O arquivo \texttt{pokemon\_names.txt} deve estar no mesmo diretório do executável para que o carregamento automático dos Pokémons funcione corretamente.

% -------------------------------------------------------
% REFERÊNCIAS
% -------------------------------------------------------
\begin{thebibliography}{1}

\bibitem{szwarcfiter2010}
SZWARCFITER, J. L.; MARKENZON, L.
\textit{Estruturas de Dados e seus Algoritmos}.
3. ed. Rio de Janeiro: LTC, 2010.

\bibitem{costa2026}
COSTA, Anderson Bessa da.
\textit{Trabalho Prático -- Disciplina de Estruturas de Dados}.
Faculdade de Computação (FACOM), Universidade Federal de Mato Grosso do Sul (UFMS), 2026.

\end{thebibliography}

\end{document}
